% Generated by scripts/assemble.py; corrections belong in source/corrections.json.
\examhead{152}{中山大学 2024 年数学分析试卷}{144}
% p145-b17; source PDF page 145
\blockanchor{p145-b17}
\question{1}(20 分) 设 $f(x) = \begin{cases} \frac{1}{x} e^{-\frac{1}{x^2}}, & x \neq 0; \\ a, & x = 0. \end{cases}$\par
% p145-b18; source PDF page 145
\blockanchor{p145-b18}
(1) 若 $f(x)$ 在 $x = 0$ 连续, 求 $a$ 的值.\par
% p145-b19; source PDF page 145
\blockanchor{p145-b19}
(2) 在第 (1) 问的条件下, 求 $f'(x)$.\par
% p145-b20; source PDF page 145
\blockanchor{p145-b20}
(3) $f(x)$ 在 $x = 0$ 处是否可以进行泰勒展开？简要说明理由.\par
% p145-b21; source PDF page 145
\blockanchor{p145-b21}
\question{2}(20 分) 证明如下结论:\par
% p145-b22; source PDF page 145
\blockanchor{p145-b22}
(1) 已知 $0 < b < a < +\infty$，那么 $\lim_{n \to \infty} \sqrt[n]{a^n + b^n} = a$.\par
% p145-b23; source PDF page 145
\blockanchor{p145-b23}
(2) 设 $f(x)$ 在 $[a, b]$ 上连续, 那么 $\lim_{n \to \infty} \sqrt[n]{\int_{a}^{b} |f(x)|^{n} \mathrm{d}x} = \max_{x \in [a, b]} |f(x)|$.\par
% p145-b24; source PDF page 145
\par\Needspace{6\baselineskip}
\blockanchor{p145-b24}
\question{3}(10 分) 已知 $0 < a < b < +\infty, f(x)$ 在 $[a, b]$ 上可导, 求证: 存在 $\xi \in (a, b)$, 使得\par
\nopagebreak[4]
% p145-b25; source PDF page 145
\blockanchor{p145-b25}
\displayformula{\frac {a f (b) - b f (a)}{a - b} = f (\xi) - \xi f ^ {\prime} (\xi).}
% p146-b00; source PDF page 146
\blockanchor{p146-b00}
\question{4}(10 分) 已知函数列 $\{f_n(x)\}$ 在区间 $I$ 上一致收敛于 $f(x)$, 且对于每一个正整数 $n$, $f_n(x)$ 均在 $I$ 上一致连续, 证明: $f(x)$ 在 $I$ 上也一致连续.\par
% p146-b01; source PDF page 146
\blockanchor{p146-b01}
\question{5}(20 分) 证明如下结论:\par
% p146-b02; source PDF page 146
\blockanchor{p146-b02}
(1) 已知级数 $\sum_{n=0}^{\infty} u_n$ 收敛, 则 $\lim_{n \to \infty} u_n = 0$.\par
% p146-b03; source PDF page 146
\blockanchor{p146-b03}
(2) 举例若 $f(x)$ 在 $[1, +\infty)$ 上连续且 $\int_{1}^{+\infty} f(x) \, \mathrm{d}x$ 收敛, 但 $\lim_{x \to +\infty} f(x) \neq 0$.\par
% p146-b04; source PDF page 146
\blockanchor{p146-b04}
(3) 若 $f(x)$ 在 $[1, +\infty)$ 上一致连续, 加上 (2) 的条件, 证明: $\lim_{x \to +\infty} f(x) = 0$.\par
% p146-b05; source PDF page 146
\blockanchor{p146-b05}
\question{6}(10 分) 证明: $f(x, y) = xy + \frac{1}{x} + \frac{1}{y}$ 在 $0 < x, y < +\infty$ 上有最小值, 无最大值.\par
% p146-b06; source PDF page 146
\blockanchor{p146-b06}
\question{7}(10 分) 求 $y = \arctan x$ 和直线 $x = 1$ 及坐标 $x$ 轴围成曲面三角形的面积.\par
% p146-b07; source PDF page 146
\blockanchor{p146-b07}
\question{8}(10 分) 计算曲面积分 $\iint_{\Sigma} xyz \, \mathrm{d}x \, \mathrm{d}y$, 其中 $\Sigma$ 为 $x^2 + y^2 + z^2 = 1$, $x \geq 0, y \geq 0$, 方向取外侧.\par
% p146-b08; source PDF page 146
\blockanchor{p146-b08}
\question{9}(20 分) 设二元函数 $u(x, y), v(x, y)$ 在区域 $D = \{(x, y) \mid x^2 + y^2 < 1\}$ 内连续可微, 且 $\frac{\partial u}{\partial y} = \frac{\partial v}{\partial x}$.\par
% p146-b09; source PDF page 146
\blockanchor{p146-b09}
(1) 试求区域 $D$ 上的二元函数 $f(x, y)$ (用 $u, v$ 来表示) 满足 $\frac{\partial f}{\partial x} = u, \frac{\partial f}{\partial y} = v$.\par
% p146-b10; source PDF page 146
\blockanchor{p146-b10}
(2) 若 $D$ 换成 $D' = \{(x, y) \mid 0 < x^2 + y^2 < 1\}$, 那么上述的函数 $f(x, y)$ 是否仍总是存在? 请说明理由.\par
% p146-b11; source PDF page 146
\blockanchor{p146-b11}
\question{10}(20 分) 解答如下问题:\par
% p146-b12; source PDF page 146
\blockanchor{p146-b12}
(1) $\int_0^{+\infty} e^{-xy} \sin y \, \mathrm{d}y$ 在 $[0, +\infty)$ 上是否一致收敛？请说明理由.\par
% p146-b13; source PDF page 146
\blockanchor{p146-b13}
(2) $\int_0^{+\infty}\mathrm{d}x\int_0^{+\infty}|e^{-xy}\sin y|\mathrm{d}y$ 是否有限？请说明理由.\par
% p146-b14; source PDF page 146
\blockanchor{p146-b14}
(3) 证明: \displayformula{\int_{0}^{+\infty} \mathrm{d}x \int_{0}^{+\infty} e^{-xy} \sin y \, \mathrm{d}y = \int_{0}^{+\infty} \mathrm{d}x \int_{0}^{+\infty} e^{-xy} \sin x \, \mathrm{d}y.}\par
